\chapter{CRUD: manipulando os dados}
\label{cap:crud}

\textbf{CRUD} é o acrônimo das quatro operações básicas sobre dados. Todo
sistema que você escrever na vida faz essas quatro coisas.

\begin{table}[H]
\centering
\caption{As quatro operações do CRUD}
\label{tab:crud}
\footnotesize
\begin{tabular}{@{}lll@{}}
\toprule
\textbf{CRUD} & \textbf{Comando SQL} & \textbf{Categoria} \\
\midrule
\textbf{C}reate & \comando{INSERT} & DML \\
\textbf{R}ead   & \comando{SELECT} & DQL \\
\textbf{U}pdate & \comando{UPDATE} & DML \\
\textbf{D}elete & \comando{DELETE} & DML \\
\bottomrule
\end{tabular}
\end{table}

\section{Create: \comando{INSERT}}

\begin{sqlcode}
-- Uma linha
INSERT INTO curso (nome) VALUES ('Engenharia de Software');

-- Varias linhas de uma vez (mais eficiente)
INSERT INTO curso (nome) VALUES
  ('Ciencia da Computacao'),
  ('Sistemas de Informacao'),
  ('Engenharia de Producao');

INSERT INTO membro (nome, email, periodo, data_ingresso, curso_id) VALUES
  ('Ana Souza',   'ana@liga.edu.br',   7, '2024-03-11', 1),
  ('Bruno Lima',  'bruno@liga.edu.br', 3, '2025-02-20', 2),
  ('Carla Dias',  'carla@liga.edu.br', 5, '2024-08-05', 1),
  ('Diego Rocha', 'diego@liga.edu.br', 9, '2023-09-14', 3),
  ('Elisa Prado', 'elisa@liga.edu.br', 1, '2026-03-02', NULL);
\end{sqlcode}

Repare que \comando{id}, \comando{ativo} e \comando{criado\_em} não foram
informados: são preenchidos pelo \comando{AUTO\_INCREMENT} e pelos
\comando{DEFAULT}.

\begin{itemize}
    \item Datas vão sempre no formato \comando{'AAAA-MM-DD'} e entre aspas.
    \item Texto entre aspas simples (\comando{'Ana'}); números sem aspas (\comando{7}).
    \item A ordem dos valores em \comando{VALUES} deve corresponder exatamente
          à ordem das colunas listadas.
\end{itemize}

\section{Read: \comando{SELECT}}

O comando mais usado da linguagem. A estrutura completa é:

\begin{sqlcode}
SELECT   colunas
FROM     tabela
WHERE    condicao          -- filtra LINHAS
ORDER BY coluna [ASC|DESC] -- ordena
LIMIT    n;                -- limita a quantidade
\end{sqlcode}

\begin{sqlcode}
SELECT * FROM membro;                    -- tudo (bom para explorar, ruim em producao)
SELECT nome, email FROM membro;          -- so as colunas que interessam
SELECT nome AS aluno, periodo AS sem FROM membro;  -- AS renomeia no resultado
\end{sqlcode}

\subsection{Filtrando com \comando{WHERE}}

\begin{sqlcode}
SELECT * FROM membro WHERE periodo > 5;
SELECT * FROM membro WHERE ativo = TRUE AND periodo >= 5;
SELECT * FROM membro WHERE periodo IN (1, 3, 5);
SELECT * FROM membro WHERE periodo BETWEEN 3 AND 7;    -- inclusivo nas 2 pontas
SELECT * FROM membro WHERE nome LIKE 'A%';             -- comeca com A
SELECT * FROM membro WHERE email LIKE '%@liga.edu.br'; -- termina com
SELECT * FROM membro WHERE nome LIKE '%ana%';          -- contem
SELECT * FROM membro WHERE curso_id IS NULL;           -- sem curso informado
SELECT * FROM membro WHERE data_ingresso >= '2025-01-01';
\end{sqlcode}

\begin{table}[H]
\centering
\caption{Operadores de filtro}
\label{tab:operadores}
\footnotesize
\begin{tabularx}{\textwidth}{@{}lX@{}}
\toprule
\textbf{Operador} & \textbf{Significado} \\
\midrule
\comando{=} \quad \comando{<>} (ou \comando{!=}) & Igual / diferente \\
\comando{>} \quad \comando{<} \quad \comando{>=} \quad \comando{<=} & Comparação \\
\comando{AND} \quad \comando{OR} \quad \comando{NOT} & Combinação lógica \\
\comando{IN (a, b, c)}      & Está na lista \\
\comando{BETWEEN a AND b}   & Está no intervalo, incluindo as pontas \\
\comando{LIKE}              & Padrão de texto: \comando{\%} é qualquer coisa, \comando{\_} é um caractere \\
\comando{IS NULL} / \comando{IS NOT NULL} & Testa ausência de valor \\
\bottomrule
\end{tabularx}
\end{table}

\subsection{Ordenando e limitando}

\begin{sqlcode}
SELECT nome, periodo FROM membro ORDER BY periodo DESC;
SELECT nome, periodo FROM membro ORDER BY periodo DESC, nome ASC;
SELECT nome FROM membro ORDER BY data_ingresso ASC LIMIT 3;  -- os 3 mais antigos
\end{sqlcode}

\comando{ASC} é crescente (padrão, pode ser omitido) e \comando{DESC} é
decrescente.

\section{Update: alterando}

\begin{sqlcode}
UPDATE membro
SET    periodo = 8
WHERE  id = 1;

UPDATE membro
SET    ativo = FALSE, periodo = 10
WHERE  email = 'diego@liga.edu.br';
\end{sqlcode}

\section{Delete: removendo}

Para praticar sem destruir os dados da aula, crie um registro descartável e
apague apenas ele:

\begin{sqlcode}
INSERT INTO membro (nome, email, periodo, data_ingresso)
VALUES ('Registro Teste', 'teste@liga.edu.br', 1, '2026-01-01');

SELECT * FROM membro WHERE email = 'teste@liga.edu.br';   -- confira ANTES
DELETE FROM membro WHERE email = 'teste@liga.edu.br';
\end{sqlcode}

\section{O erro mais caro de quem está começando}
\label{sec:where-ausente}

\begin{aviso}[Sem \comando{WHERE}, o comando vale para a tabela inteira]
\begin{sqlcode}
UPDATE membro SET periodo = 1;   -- zera o periodo de TODOS os membros
DELETE FROM membro;              -- apaga TODOS os membros da tabela
\end{sqlcode}
E não existe ``desfazer'' em SQL.
\end{aviso}

O hábito obrigatório é testar com \comando{SELECT} antes:

\begin{sqlcode}
-- 1) veja EXATAMENTE quais linhas serao afetadas
SELECT * FROM membro WHERE id = 5;

-- 2) so entao troque o SELECT * pelo DELETE ou UPDATE,
--    mantendo exatamente o mesmo WHERE
DELETE FROM membro WHERE id = 5;
\end{sqlcode}

\newpage

\begin{dica}[Rede de proteção, e a falta dela]
O MySQL Workbench tem o \emph{safe update mode}, que \textbf{bloqueia}
\comando{UPDATE} e \comando{DELETE} sem \comando{WHERE} em coluna-chave. O
phpMyAdmin \textbf{não} protege: nele, a disciplina é sua.

\comando{DELETE FROM tabela;} apaga as linhas uma a uma (é DML).
\comando{TRUNCATE TABLE tabela;} esvazia a tabela de forma muito mais rápida e
reinicia o \comando{AUTO\_INCREMENT}, mas é DDL e não pode ser desfeito por
\comando{ROLLBACK}.
\end{dica}
